\section{A threshold for one negation}
Two monotone restrictions give a lower bound for circuits with one NOT gate.

Write \mbox{$(a-1)_+=\max\{a-1,0\}$} for a nonnegative integer $a$.
Let $S_1,\ldots,S_m$ be subsets of a finite input set $V$, with distinct sets meeting in at most one point. Their designated outputs are the disjunctions $\bigvee_{v\in S_i}x_v$.
The following theorem is conditional on the stated monotone lower bound; this separates the restriction argument from its base case.

\begin{theorem}\label{thm:one-not-threshold}
Assume that every family of subsets of any finite set with pairwise intersection at most one requires at least
$\sum_i(|S_i|-1)_+$ gates to compute all its disjunctions by a monotone circuit of fan-in at most two.
Let $C$ be an AND/OR/NOT circuit of fan-in at most two computing the disjunctions of a family $S_1,\ldots,S_m\subseteq V$ with this intersection property. Suppose $C$ has at most one NOT gate, and put $n=|V|$.
There is an integer $W$, $0\leq W\leq n+1$, such that
\begin{enumerate}
\item if $W\geq1$, some $a\subseteq V$ of size $W-1$ satisfies
\[
 \sum_i(|S_i\cap a|-1)_+\leq |C|;
\]
\item for every $b\subseteq V$ of size $W$,
\[
 \sum_{i:S_i\cap b=\varnothing}(|S_i|-1)_+\leq |C|.
\]
\end{enumerate}
For $W=n+1$ the second condition is vacuous.
\end{theorem}
\begin{proof}
If there is no NOT gate, take $W=0$ and apply the assumed bound to the whole family. Otherwise let $h(x)$ be the Boolean value on the input wire of the unique NOT gate. It is monotone, because its preceding circuit has no negation.
If $h$ is identically one, put $W=0$. Otherwise put
\[
 W=1+\max\{|a|:a\subseteq V,\ h(\mathbf1_a)=0\}.
\]
Choose a maximizing $a$. Monotonicity gives $h(x)=0$ on the down-set $x\leq\mathbf1_a$. The definition gives $h(\mathbf1_b)=1$ for every set $b$ of size $W$, and then $h(x)=1$ on the up-set $x\geq\mathbf1_b$. If $h$ is identically zero, this construction gives $W=n+1$.

Fix all inputs outside $a$ to zero. The NOT output becomes constant one, and the required outputs become the disjunctions of $S_i\cap a$. Propagate constants in topological order: a gate with a forcing constant becomes a constant; a gate with an identity constant redirects to its other operand; a gate with two live operands is kept. After removing designated outputs of size zero or one, every remaining output is nonconstant. Thus no constant-output gate needs to be introduced, and the resulting monotone circuit has at most $|C|$ gates. The base hypothesis applies, since intersections cannot increase. Discarded sets contribute zero to the asserted sum. This proves the first inequality.

For a fixed $b$ of size $W$, set its inputs to one and keep only the outputs with $S_i\cap b=\varnothing$. The NOT output is constant zero. The same propagation gives a monotone circuit of no greater size for the disjunctions of these sets on $V\setminus b$. Again discard sets of size at most one and apply the base hypothesis. This proves the second inequality, including $W=0$ when $h$ is identically one.
\end{proof}

\begin{corollary}\label{cor:threshold-bound}
Under the theorem's hypothesis, if a number $K$ has the following property for every $W\in\{0,\ldots,n+1\}$, then $|C|\geq K$: either $W\geq1$ and every $(W-1)$-element set $a$ has first sum at least $K$, or some $W$-element set $b$ has second sum at least $K$.
\end{corollary}
\begin{proof}
Take the value $W$ from Theorem~\ref{thm:one-not-threshold}. Its witnessing set $a$ gives the conclusion in the first alternative; its inequality for every $b$ gives the conclusion in the second.
\end{proof}

For a projective plane of order $q$ there are $N=q^2+q+1$ lines, each of size $q+1$. Let $e(A)$ be the number of lines disjoint from a point set $A$. Incidence counting gives
\[
 A_k:=\min_{|a|=k}\sum_{\ell}(|\ell\cap a|-1)_+
       =(q+1)k-N+\min_{|a|=k}e(a),
 \qquad B_k:=q\max_{|b|=k}e(b).
\]
In the threshold expression take $A_{-1}=B_{N+1}=0$.

\begin{proposition}\label{prop:plane-four-threshold}
For the Desarguesian projective plane of order four,
\[
 \min_{0\leq W\leq22}\max\{A_{W-1},B_W\}=14.
\]
Thus the preceding restriction principle yields a fourteen-gate lower bound for its line disjunctions with one NOT gate, conditional on its monotone base hypothesis. The number fourteen is the exact value of this threshold expression; it is not asserted to be the optimum circuit size or the strongest available lower bound.
\end{proposition}
\begin{proof}
There are twenty-one points and lines, with five incidences at each point. Four lines with no three concurrent have union of size $4\cdot5-\binom42=14$. Their seven-point complement avoids all four lines. Every subset of this complement also does so; hence $B_W\geq16$ for $W\leq7$.
For $W\geq8$, incidence counting and $e(a)\geq0$ give
$A_{W-1}\geq5(W-1)-21\geq14$. These two observations prove the lower bound fourteen for the minimum.

At $W=8$, the seven points of the embedded plane $\mathbb P^2(\mathbb F_2)\subseteq\mathbb P^2(\mathbb F_4)$ meet every line. Indeed a linear form over $\mathbb F_4$ restricts to an $\mathbb F_2$-linear map $\mathbb F_2^3\to\mathbb F_4$, whose kernel has a nonzero vector. Thus $A_7=35-21=14$.
Any four distinct lines have union of size at least fourteen: when adding them one at a time, the new line loses at most its number of predecessors in intersections. If an eight-point set avoided four lines, their union would lie in its thirteen-point complement, a contradiction. Therefore $B_8\leq12$. Three nonconcurrent lines have union of size twelve, so an eight-point subset of their nine-point complement shows $B_8=12$. Consequently the value at $W=8$ is fourteen, proving equality.
\end{proof}

A separate exhaustive calculation of the finite-plane subset calculation gives threshold values $2$, $6$ and $14$ for orders $2$, $3$ and $4$, respectively. At order four it checks all $2^{21}$ point subsets and finds $168$ six-point hyperovals. The geometric proof above establishes the value fourteen without relying on this enumeration.
